\title[Reliability Boundaries of HyDE-Style Query Expansion]{Reliability Boundaries of HyDE-Style Query Expansion in Multi-Hop Intelligent Information Retrieval}

\author*[1]{\fnm{Shiyong} \sur{Xiong}}\email{1771856981@qq.com}
\author[1]{\fnm{Xingyun} \sur{Chen}}

\affil[1]{\orgdiv{School of Computer Science and Technology}, \orgname{Chongqing University of Posts and Telecommunications}, \orgaddress{\city{Chongqing}, \postcode{400065}, \country{China}}}

\abstract{Generative query expansion can improve retrieval without establishing corresponding gains in evidence delivered to readers or in multi-hop answer quality. We examine these configuration-specific reliability boundaries for Hypothetical Document Embeddings (HyDE)-style query expansion, using matched question-only baselines, complete-support retrieval, actual-input audits, and paired answer comparisons with family-specific Holm correction. In a frozen hosted HotpotQA artifact, HyDE raises all-support hit@10 from 0.630 to 0.868, but 292 of 500 generated passages contain exact answers; in separate matched BGE-encoder comparisons, two pinned 7B generators do not reproduce the hosted positive direction. In a random 300-question HotpotQA comparison excluding previously evaluated IDs and using a 5.23-million-paragraph corpus, cross-encoder reranking increases complete-support retrieval from 0.350 to 0.537 and exact-match accuracy by 9.0 percentage points with the Qwen reader and 7.0 with the Mistral reader, relative to BM25. HyDE improves retrieval, but its answer gain over BM25 remains unresolved with the Mistral reader. Cross-encoder retrieval gains on the 2WikiMultihopQA split-union corpus likewise do not establish answer gains. On MultiHop-RAG news documents, standard HyDE underperforms the matched question-only dense baseline in complete-support retrieval and all eight reader/length exact-match comparisons. In a separate matched 100-question HotpotQA intervention, a larger character budget increases complete-support retention from 46\% to 55\% without establishing answer gains. These findings distinguish retrieval, delivered evidence, and answers as separate validation targets within the tested configurations. The absence of an iterative-retrieval comparator on the external news corpus, incomplete cost measurements, and the absence of independent human validation constrain generalization.}
\keywords{intelligent information retrieval; retrieval-augmented generation; multi-hop question answering; query expansion; reproducibility}

\maketitle
